% Archivo: w12_s03_brief_bom_aprobada.tex
\documentclass[11pt, a4paper]{article}

\usepackage[T1]{fontenc}
\usepackage[utf8]{inputenc}
\usepackage{tgpagella}
\usepackage[spanish, es-noshorthands]{babel}
\usepackage{amsmath, amssymb, bm}
\usepackage{tabularx}
\usepackage{booktabs}
\usepackage[most]{tcolorbox}
\usepackage[margin=2.2cm]{geometry}
\usepackage{ragged2e}
\usepackage{fancyhdr}
\usepackage{tikz}
\usepackage[american]{circuitikz}

\pagestyle{fancy}
\fancyhf{}
\setlength{\headheight}{16pt}
\lhead{\textbf{UCB Sede Tarija} | Ing. Mecatrónica}
\chead{}
\rhead{IMT-121 Circuitos Electrónicos I | W12S3 — Artefacto A3}
\lfoot{Prof: M.Sc. Ing. Bernardo Quiroga Turdera}
\rfoot{Página \thepage}
\renewcommand{\headrulewidth}{0.4pt}

\definecolor{ucbblue}{RGB}{0, 51, 102}
\definecolor{warnred}{RGB}{180, 0, 0}

\begin{document}

\begin{tcolorbox}[colback=ucbblue!5!white, colframe=ucbblue, title=\textbf{Artefacto A3: Brief Técnico RC–RLC y BOM de Compra Aprobada (W12S3)}]
    \centering \textbf{Preparación de Experiencia Integrada (26/10 y 28/10) — Compra Aprobada $\neq$ Autorización de Energizar}
\end{tcolorbox}

\section*{1. Diseño Candidato de Dos Etapas con Generador de Señales (T3 y Visual 4)}
Para unificar materiales y utilizar un único capacitor no polarizado de $C = 1\,\mu\text{F}$ ($\ge 25\,\text{V}$), se sustituye la variante histórica de batería de $9\,\text{V}$ y capacitor electrolítico de $470\,\mu\text{F}$. Ambas etapas se excitan con generador de funciones (onda cuadrada/escalón cuyos niveles y frecuencias definitivos se fijarán en la prueba piloto docente).

\vspace{0.25cm}
\noindent
\begin{minipage}[t]{0.485\textwidth}
\vspace{0pt}
\begin{tcolorbox}[colback=white, colframe=ucbblue, left=2mm, right=2mm, top=2mm, bottom=2mm, title={\small\textbf{Etapa A: Circuito RC Candidato}}]
\centering
\begin{circuitikz}[scale=0.72, transform shape, thick]
    \draw (0,0) to[sV, l_={Gen. $v_s(t)$}] (0,2.4)
          to[short, -*] (0.9,2.4) node[above]{\texttt{in}}
          to[R, l={$R_{\text{ext}}=10\,\text{k}\Omega$}] (3.6,2.4)
          to[short, -*] (4.3,2.4) node[above]{\texttt{out}};
    \draw (4.3,2.4) to[C, l={$C=1\,\mu\text{F}$}, v={$v_C$}] (4.3,0);
    \draw[ucbblue] (0.9,2.4) to[short, *-o] (0.9,3.2) node[above]{\scriptsize CH1: $V(\text{in})$};
    \draw[ucbblue] (4.3,2.4) to[short, *-o] (5.1,2.4) node[right]{\scriptsize CH2: $V(\text{out})$};
    \draw (0,0) -- (5.1,0) node[right]{\scriptsize Retorno};
    \draw (2.2,0) node[ground]{};
\end{circuitikz}
\par\vspace{0.15cm}
\RaggedRight\footnotesize
\textbf{Predicción nominal:} $R_{\text{eq}} \approx R_{\text{ext}} + R_{\text{gen}}$. Despreciando provisionalmente $R_{\text{gen}}$:
$$\tau_{\text{nom}} \approx (10\,000\,\Omega)(10^{-6}\,\text{F}) = \mathbf{10\,\text{ms}}$$
$$v_C(s) = V_f + (V_i - V_f)e^{-s/\tau}$$
\textbf{Ejemplo didáctico WE3:} Para un escalón hipotético $0 \to 5\,\text{V}$, en $s=\tau$ se alcanza $5(1-e^{-1}) \approx 3.16\,\text{V}$ (para $0 \to 10\,\text{V}$, $v_C(\tau) \approx 6.32\,\text{V}$; amplitudes sujetas a piloto).
\end{tcolorbox}
\end{minipage}%
\hfill
\begin{minipage}[t]{0.485\textwidth}
\vspace{0pt}
\begin{tcolorbox}[colback=white, colframe=ucbblue, left=2mm, right=2mm, top=2mm, bottom=2mm, title={\small\textbf{Etapa B: Circuito RLC Serie Candidato}}]
\centering
\begin{circuitikz}[scale=0.68, transform shape, thick]
    \draw (0,0) to[sV, l_={Gen. $v_s(t)$}] (0,2.4)
          to[short, -*] (0.8,2.4) node[above]{\texttt{in}}
          to[R, l={$R_{\text{ext}}$}] (2.7,2.4)
          to[L, l={$L=10\,\text{mH}$}] (4.7,2.4)
          to[short, -*] (5.3,2.4) node[above]{\texttt{out}};
    \draw (5.3,2.4) to[C, l={$C=1\,\mu\text{F}$}, v={$v_C$}] (5.3,0);
    \draw[ucbblue] (0.8,2.4) to[short, *-o] (0.8,3.2) node[above]{\scriptsize CH1: $V(\text{in})$};
    \draw[ucbblue] (5.3,2.4) to[short, *-o] (6.0,2.4) node[right]{\scriptsize CH2: $V(\text{out})$};
    \draw (0,0) -- (6.0,0) node[right]{\scriptsize Retorno};
    \draw (2.7,0) node[ground]{};
\end{circuitikz}
\par\vspace{0.15cm}
\RaggedRight\footnotesize
\textbf{Predicción nominal:} $\omega_0 \approx 10\,000\,\text{rad/s}$ ($f_0 \approx 1.59\,\text{kHz}$, $T_0 \approx 0.628\,\text{ms}$), $R_{\text{crit,total}} \approx 200\,\Omega$.
$$R_{\text{total}} = R_{\text{ext}} + R_{\text{gen}} + R_{\text{DCR}} + R_{\text{otros}}$$
$$R_{\text{ext,crit}} \approx 200\,\Omega - R_{\text{gen}} - R_{\text{DCR}} - R_{\text{otros}}$$
\textbf{Exploración BOM:} Se emplean $R_{\text{ext}} \in \{47, 100, 330\}\,\Omega$ para contrastar el amortiguamiento real tras medir $R_{\text{DCR}}$ y verificar $R_{\text{gen}}$.
\end{tcolorbox}
\end{minipage}

\section*{2. Visual 3 — Comparación de Escalas Temporales (Procedimental)}
No se debe configurar una única ventana horizontal en el osciloscopio para ambas etapas:

\vspace{0.15cm}
\begin{center}
\begin{tikzpicture}[scale=0.92, transform shape]
    % Escala RC
    \node[anchor=east, font=\bfseries\footnotesize, ucbblue] at (-0.2, 1.3) {Etapa A (RC):};
    \draw[thick, -Stealth, ucbblue] (0,1.3) -- (8.8,1.3) node[right]{\footnotesize Escala: decenas de ms};
    \foreach \x/\lbl in {0/0 ms, 1.8/$\tau \approx 10$ ms (63.2\%), 7.6/$5\tau \approx 50$ ms} {
        \draw[thick, ucbblue] (\x,1.15) -- (\x,1.45);
        \node[above, font=\scriptsize] at (\x,1.45) {\lbl};
    }
    % Escala RLC
    \node[anchor=east, font=\bfseries\footnotesize, warnred] at (-0.2, 0) {Etapa B (RLC):};
    \draw[thick, -Stealth, warnred] (0,0) -- (8.8,0) node[right]{\footnotesize Escala: pocos ms ($\mu$s/muestra)};
    \foreach \x/\lbl in {0/0 ms, 2.4/$T_0 \approx 0.628$ ms (1 ciclo), 7.2/Ventana $\sim 2\text{--}3$ ms} {
        \draw[thick, warnred] (\x,-0.15) -- (\x,0.15);
        \node[below, font=\scriptsize] at (\x,-0.15) {\lbl};
    }
\end{tikzpicture}
\end{center}

\section*{3. Tabla BOM con Compra Aprobada y Trazabilidad de Estados (T4)}
\noindent\textit{Distinción obligatoria de estados: \textbf{(1) Compra aprobada} $\to$ \textbf{(2) Adquirido} $\to$ \textbf{(3) Verificado con multímetro} $\to$ \textbf{(4) Apto para montaje}. Reserva docente sugerida (no obligatoria): 11 inductores y 12 unidades de cada resistor/capacitor.}

\vspace{0.2cm}
{\small
\setlength{\tabcolsep}{3.5pt}
\renewcommand{\arraystretch}{1.45}
\noindent
\begin{tabularx}{\textwidth}{
  >{\bfseries\RaggedRight\hsize=0.85\hsize}X
  >{\centering\arraybackslash\hsize=0.65\hsize}X
  >{\RaggedRight\arraybackslash\hsize=1.40\hsize}X
  >{\centering\arraybackslash\hsize=0.65\hsize}X
  >{\centering\arraybackslash\hsize=1.00\hsize}X
  >{\RaggedRight\arraybackslash\hsize=1.45\hsize}X}
\toprule
\textbf{Elemento} & \textbf{Cant.} & \textbf{Característica Aprobada} & \textbf{Total*} & \textbf{Estado Normativo} & \textbf{Verificación Física Grupo} \\
\midrule
Inductor $L$    & 1 & $10\,\text{mH}$, pines aptos para protoboard; medir DCR/corriente & 9 & Compra aprobada & {[$\square$]} Adq. {[$\square$]} Verif. $R_{\text{DCR}}$:\,\rule{0.7cm}{0.4pt} \\
Capacitor $C$   & 1 & $1\,\mu\text{F}$ \textbf{no polarizado} (película pref.), $\ge 25\,\text{V}$ & 9 & Compra aprobada & {[$\square$]} Adq. {[$\square$]} Verif. {[$\square$]} Apto \\
Resistor RC     & 1 & $10\,\text{k}\Omega$, $1/4\,\text{W}$, $\le 5\%$ & 9 & Compra aprobada & {[$\square$]} Adq. {[$\square$]} Medido: \rule{0.9cm}{0.4pt} \\
Resistor bajo   & 1 & $47\,\Omega$, $1/4\,\text{W}$, $\le 5\%$  & 9 & Compra aprobada & {[$\square$]} Adq. {[$\square$]} Medido: \rule{0.9cm}{0.4pt} \\
Resistor medio  & 1 & $100\,\Omega$, $1/4\,\text{W}$, $\le 5\%$ & 9 & Compra aprobada & {[$\square$]} Adq. {[$\square$]} Medido: \rule{0.9cm}{0.4pt} \\
Resistor alto   & 1 & $330\,\Omega$, $1/4\,\text{W}$, $\le 5\%$ & 9 & Compra aprobada & {[$\square$]} Adq. {[$\square$]} Medido: \rule{0.9cm}{0.4pt} \\
Prototipado     & 1 + jgo. & Protoboard y jumpers (si lab no provee) & 9 + jgos. & Por confirmar & {[$\square$]} Disponible {[$\square$]} Revisado \\
\bottomrule
\end{tabularx}
}
\par\vspace{0.1cm}{\footnotesize *Total para 9 grupos condicionado a la confirmación definitiva del número real de grupos.}

\vspace{0.2cm}
\begin{tcolorbox}[colback=gray!5, colframe=black, title=\textbf{Registro de Alternativas Aprobadas Antes de Energizar (Si el componente adquirido difiere)}]
Valor real / alternativa propuesta: \hrulefill \par\vspace{0.15cm}
Recálculo de $\tau$ / $\omega_0$ / $R_{\text{crit}}$: \rule{6.2cm}{0.4pt} \hfill Firma V.º B.º Docente: \rule{3.2cm}{0.4pt}
\end{tcolorbox}

\begin{tcolorbox}[colback=warnred!5!white, colframe=warnred, title=\textbf{Comunicado Oficial para Estudiantes (Emisión Docente)}]
\small\RaggedRight
Para EC3 traer el lunes 26/10 el kit RC–RLC aprobado por grupo: una bobina de $10\,\text{mH}$ apta para protoboard, un capacitor no polarizado de $1\,\mu\text{F}$ ($\ge 25\,\text{V}$), resistencias de $10\,\text{k}\Omega$, $47\,\Omega$, $100\,\Omega$ y $330\,\Omega$ ($1/4\,\text{W}$, $\le 5\%$), y protoboard/cables si el laboratorio no los proporciona. El trabajo experimental del 26/10 y 28/10 depende de verificación y autorización de equipos y seguridad; traer materiales \textbf{no} implica autorización para energizar. El examen individual está previsto para el 30/10. La entrega digital del informe IEEE el 06/11 a las 23:59 continúa como plazo propuesto pendiente de confirmación formal.
\end{tcolorbox}

\end{document}
